% ===================================================================
%  اللوحة رقم ٥ — البيت وبناه.
%
%  Traduction, faite en 2026, en arabe levantin d'aujourd'hui, sur
%  l'ido de texto/io. Regles de la colonne : voir 10-lawha-01.tex.
%  Tableau en DIALOGUE d'un bout a l'autre — trente-six attributions
%  de parole —, sans un seul alinea numerote par l'auteur : toutes
%  les cles sont donc « t05-01-n ». Deux renvois portent « bis »,
%  94 bis et 95 bis, recopies tels quels.
%
%  LE CHANTIER EST TURC DU SOL AU TOIT, ET C'EST LA SUITE EXACTE DU
%  TABLEAU 1. Le poele et son tuyau y etaient turcs quand l'ecole
%  restait arabe : les Ottomans avaient meuble la piece. Ici ils
%  nomment les METIERS. Le charretier (135) est « العربجي »,
%  l'electricien (124) « الكهربجي », le couvreur (118) « القرميدجي » :
%  c'est le suffixe turc -ci, colle sur des racines arabes, et il est
%  encore productif — un Levantin fabrique un nom de metier avec
%  n'importe quel objet en lui ajoutant جي, et le fait tous les jours.
%  QUATRE SIECLES D'ADMINISTRATION N'ONT PAS LAISSE UNE LANGUE, ILS
%  ONT LAISSE UNE TERMINAISON.
%
%  ET LES MATERIAUX SONT ITALIENS. Le mortier (65) est « المونة »,
%  la truelle « المالج », l'echafaudage (110) « السقالة » : les
%  batisseurs genois et venitiens ont tenu les ports de cette cote
%  pendant des siecles, et ce qu'ils ont laisse est exactement ce
%  qu'on melange et ce sur quoi on monte. Le metier est turc, la
%  matiere est italienne, ET LA MAISON EST ARABE — الحوش la cour
%  (45), الحمّام le bain, العلّية le grenier, العتبة le seuil (10).
%
%  L'ALCOVE (h) EST « الليوان », ET LES DEUX MOTS ONT FAIT LE MEME
%  VOYAGE EN SENS INVERSE. Le liwan est la piece voutee ouverte sur
%  la cour de la maison damascene, et le mot est persan. « Alcove »,
%  lui, est l'arabe القبّة parti en Espagne et revenu par le francais.
%  Deux mots pour un renfoncement, chacun emprunte a la direction
%  d'ou l'autre vient.
%
%  Enfin le volet (19) est « الشيش ». Le mot avait servi au tableau 2
%  pour le fleuret de l'escrimeur — meme mot, deux objets — et il a
%  ete rendu ici a son sens premier : le tableau 2 ecrit desormais
%  « سيف المبارزة ». UNE COLONNE PEUT SE CORRIGER EN ARRIERE, tant
%  qu'elle le fait avant que le mot serve deux fois.
% ===================================================================

%%K t05-apar-1 apar
\VUsaut{6.0mm}
\VUtitre{15.0pt}{60}{اللوحة رقم ٥}
\VUsaut{1.4mm}
\VUtitre{11.4pt}{0}{البيت وبناه.}
\VUsaut{2.2mm}

%%K t05-tit-1 sub
\VUpk{10.2pt}{40}{لقا منيح.}
\VUsaut{3.0mm}

%%K t05-01-1 p
\textsc{حنّا}. --- يا عمّي، كتير بدّي أعرف كيف بيبنوا \VUgras{بيت}.

%%K t05-01-2 p
\textsc{العمّ} (80). --- هاد شي كتير سهل يا حبيبي، تعا معي، رح فرجيك
هداك اللي السيّد برناردوس (79) عم يبنيه واللي أنا رح سكن فيه.

%%K t05-01-3 p
\textsc{حنّا}. --- ومين رسم \VUgras{مخطّط} \textsuperscript{(76)} هالبيت ؟

%%K t05-01-4 p
\textsc{العمّ}. --- \VUgras{المهندس المعماري} \textsuperscript{(75)} ؛ قدّم رسم
كتير واضح وانقبل، و\VUgras{المتعهّد} \textsuperscript{(77)} بلّش الشغل.

%%K t05-01-5 p
\textsc{حنّا}. --- ومين هوّي المتعهّد ؟

%%K t05-01-6 p
\textsc{العمّ}. --- هوّي السيّد بلانكو، أبو رفيقك يعقوب. عم يدير
العمّال مع السيّد روفو، المهندس المعماري.

%%K t05-01-7 p
هيّه ! هيّن جايين بالوقت السيّد بلانكو مع \VUgras{مسطرته}
\textsuperscript{(78)}، والسيّد روفو مع مخطّطه.

%%K t05-01-8 p
صباح الخير يا سادة. كيفكن ؟

%%K t05-01-9 p
\textsc{السيّد روفو}. --- منيحين، شكراً. صباح الخير يا حنّا، شو ما
كبر هالولد !

%%K t05-01-10 p
\textsc{العمّ}. --- إي والله، صار رجّال زغير. كتير جدّي، لازم تشرحلو كل
شي بيشوفه. اليوم بدّه يعرف كيف بيبنوا بيت.

%%K t05-tit-2 sub
\VUpk{10.2pt}{40}{العمّال.}
\VUsaut{3.0mm}

%%K t05-01-11 p
\textsc{السيّد بلانكو}. --- طيّب، تعا معي يا حبيبي، رح شرحلك حالاً،
لأني كتير بحبّ الولاد اللي بدّهن يتعلّموا.

%%K t05-01-12 p
عم تشوف هالعامل اللي ماسك \VUgras{رفش} \textsuperscript{(82)} بإيده : هوّي
\VUgras{عامل الحفر} \textsuperscript{(81)}. عم يحفر \VUgras{خندق}
\textsuperscript{(46)} ليحطّ مواسير الميّ. وكمان حفر الأرض بالمحل اللي
فيه البيت، ليقدر \VUgras{البنّا} يبني الحيطان الأربعة. والقسم من
هالحيطان اللي بالأرض بيسمّوه \VUgras{الأساسات} \textsuperscript{(2)}. والفراغ
اللي بين الأساسات بيعمل \VUgras{القبو} \textsuperscript{(3)}. وهالقبو إلو
شبّاك زغير بيسمّوه \VUgras{طاقة القبو} \textsuperscript{(5)}، و\VUgras{باب}
\textsuperscript{(4)} ودرج.

%%K t05-01-13 p
و\VUgras{البنّا} \textsuperscript{(108)} كمّل يبني الحيطان وترك فتحات
لـ\VUgras{الأبواب} \textsuperscript{(8)} و\VUgras{الشبابيك} \textsuperscript{(17)}.

%%K t05-01-14 p
\textsc{حنّا}. --- بس أنا ما عم شوف هالبنّا !

%%K t05-01-15 p
\textsc{السيّد بلانكو}. --- هلّق خلّص شغله بالبيت. هوّي جنب
\VUgras{الإسطبل} \textsuperscript{(54)}، على \VUgras{سقالته} \textsuperscript{(110)}،
عم يبني حيط \VUgras{المغسلة} \textsuperscript{(56)}.

%%K t05-01-16 p
معه مالج وطشت و\VUgras{شاقول} \textsuperscript{(109)}. بس إنت ما رح تتذكّر
هالأسامي.

%%K t05-01-17 p
\textsc{حنّا}. --- أكيد رح تذكّرهن، لأني حالاً رح طلب من عمّي مالج
زغير وطشت، وأنا بنفسي رح عمل شاقول بخيط.

%%K t05-tit-3 sub
\VUsaut{3.0mm}
\VUpk{10.2pt}{40}{مواد البنا.}
\VUsaut{3.0mm}

%%K t05-01-18 p
\textsc{العمّ}. --- آه ! كتير منيح. بس قلّي يا حنّا، شو بدّهن ليبنوا
بيت ؟

%%K t05-01-19 p
\textsc{حنّا}. --- بدّهن \VUgras{حجار مقصوبة} \textsuperscript{(59)}
و\VUgras{مونة} \textsuperscript{(65)}.

%%K t05-01-20 p
\textsc{العمّ}. --- مظبوط. شوف هالرجّال، مع برنيطته الكبيرة، على
\VUgras{عربايته} \textsuperscript{(136)}. هوّي \VUgras{العربجي}
\textsuperscript{(135)}. كل يوم بيجيب لهون \VUgras{بلوكات حجر}
\textsuperscript{(60)}. بيفرّغ عربايته بالحوش، و\VUgras{النحّاتين}
\textsuperscript{(83)} بينحتوا البلوكات وبينشروها بـ\VUgras{منشار الحجر}
\textsuperscript{(85)}. و\VUgras{العتّال} \textsuperscript{(146)} بيحملها للبنّا
اللي بيركّبها.

%%K t05-01-21 p
وبنفس الوقت، \VUgras{عجّان المونة} \textsuperscript{(87)} عم يحضّر المونة.
بدّه \VUgras{رمل} \textsuperscript{(62)} و\VUgras{كلس} \textsuperscript{(63)}
و\VUgras{ميّ} \textsuperscript{(64)}. والكلس جنبه بـ\VUgras{كيس}
\textsuperscript{(71)}، و\VUgras{مساعد البنّا} \textsuperscript{(111)} عم
يجيبلو الرمل على \VUgras{عربايّة اليد} \textsuperscript{(112)}،
و\VUgras{الصبي} \textsuperscript{(113)} عند المضخّة (58). عم يعبّي
\VUgras{سطل} \textsuperscript{(114)}. والمونة عن قريب بتصير جاهزة.
و\VUgras{ناقل المونة} \textsuperscript{(107)} بياخدها وبيطلّعها، عالسلّم،
عالسقالة، وين عم يشتغل بنّا عم يخلّص هالجهة من البيت.

%%K t05-01-22 p
وهلّق، شو بيسمّوا العامل اللي عم يشتغل بـ\VUgras{السقف}
\textsuperscript{(31)} ؟

%%K t05-01-23 p
\textsc{حنّا}. --- هوّي \VUgras{القرميدجي} \textsuperscript{(118)}.

%%K t05-tit-4 sub
\VUsaut{3.0mm}
\VUpk{10.2pt}{40}{سقف البيت.}
\VUsaut{3.0mm}

%%K t05-01-24 p
\textsc{السيّد بلانكو}. --- إي، بس أول شي بيجي \VUgras{نجّار العمار}
\textsuperscript{(115)}.

%%K t05-01-25 p
نجّار العمار بيشتغل \VUgras{الهيكل الخشبي} \textsuperscript{(35)}. بيركّب
\VUgras{الجسور} \textsuperscript{(37)} بـ\VUgras{الأوتاد} \textsuperscript{(117)}.
وهدول العمّال، هونيك فوق، بكمامهن العريضة من مخمل، هنّي نجّارين عمار.
واحد ماسك جسر زغير، والتاني عم يقصّ وتد بـ\VUgras{فاسه}
\textsuperscript{(116)}. واللي جنب \VUgras{المدخنة} \textsuperscript{(41)} هوّي
القرميدجي. عم يسمّر شرايح على (*) العوارض، و\VUgras{أردواز} \textsuperscript{(66)} عالشرايح.
وبعض المرّات بيحطّ قرميد.

%%K t05-noto-1 noto
\VUnotes{143.2mm}{%
(*) \VUgras{Balk-o} = D. \textit{Balken}، E. \textit{joist}، F.
\textit{solive}، I. \textit{travicello}، R. \textit{balka}، S. Port.
\textit{viga}. جذر تقني ما انقبل لهلّق، بس بيستاهل يُقترح ليترجم معنى
الكلمة الفرنسية \textit{solive}، اللي مش trabo، F. \textit{poutre}، ولا
trabeto. هيّي عمود خشب مربّع بيحمل أرض الطابق والسقف.%
}

%%K t05-01-26 p
سقف الإسطبل \VUgras{مقرمد} \textsuperscript{(55)}، وسقف البيت
\VUgras{مغطّى بالأردواز} \textsuperscript{(32)}، وسقف الفرندة \VUgras{من
زنك} \textsuperscript{(24)}.

%%K t05-01-27 p
\textsc{حنّا}. --- والقرميدجي بيشتغل كمان سقوف الزنك ؟

%%K t05-01-28 p
\textsc{السيّد بلانكو}. --- لأ، بس \VUgras{عامل الزنك}
\textsuperscript{(121)} أو \VUgras{عامل الرصاص}، هوّي اللي بيحطّ
\VUgras{الكوّة} \textsuperscript{(38)} على سقف البيت. وبيحطّ كمان
\VUgras{مزاريب المطر} \textsuperscript{(43)} و\VUgras{مجاري السقف}
\textsuperscript{(42)}. بيلحّم الزنك والصفيح. وهوّي ثبّت \VUgras{مانعة
الصواعق} \textsuperscript{(40)} و\VUgras{ديك الهوا} \textsuperscript{(39)}
على \VUgras{الجملون} \textsuperscript{(36)}.

%%K t05-01-29 p
\textsc{حنّا}. --- هيك بيكون البيت خلص ؟

%%K t05-tit-5 sub
\VUsaut{3.0mm}
\VUpk{10.2pt}{40}{العدّة.}
\VUsaut{3.0mm}

%%K t05-01-30 p
\textsc{السيّد بلانكو}. --- مش كل شي. \VUgras{المبيّض}
\textsuperscript{(99)} عم يبني بـ\VUgras{الطوب} \textsuperscript{(61)}
الحيطان اللي بتقسّمه لغرف. وبيلبّس بـ\VUgras{الجصّ} \textsuperscript{(70)}
\VUgras{السقوف} \textsuperscript{(26)} والحيطان. هلّق عم يشتغل سقف
\VUgras{الفرندة} \textsuperscript{(33)}. وعنده، متل البنّا، \VUgras{مالج}
\textsuperscript{(100)} و\VUgras{طشت} \textsuperscript{(101)}.

%%K t05-01-31 p
وبعدين النجّار بيشتغل الأبواب والشبابيك و\VUgras{الباركيه}
\textsuperscript{(16)} و\VUgras{الشيش} \textsuperscript{(19)}.

%%K t05-01-32 p
هلّق عم يشتغل بالحوش على \VUgras{طاولة الشغل} \textsuperscript{(90)}.
عنده \VUgras{منشار} \textsuperscript{(94)} لينشر الخشب (69)،
و\VUgras{فارة} \textsuperscript{(91)}، و\VUgras{ملزمة} \textsuperscript{(92)}،
و\VUgras{إزميل} \textsuperscript{(94 bis)}، و\VUgras{فرجار}
\textsuperscript{(95)}، و\VUgras{شاكوش خشب} \textsuperscript{(95 bis)}.

%%K t05-01-33 p
\textsc{حنّا}. --- آه ! بس النجارة أحلى من البنا. يا عمّي، رح تشتريلي
طاولة شغل ومنشار وفارة، لأني عن قريب رح عمل بيت لميدور.

%%K t05-01-34 p
\textsc{العمّ}. --- إي يا حبيبي.

%%K t05-01-35 p
\textsc{حنّا}. --- شو ما أنا مبسوط ! ومين هاد اللي واقف جنب باب
الدخول ؟

%%K t05-tit-6 sub
\VUsaut{3.0mm}
\VUpk{10.2pt}{40}{الدهان.}
\VUsaut{3.0mm}

%%K t05-01-36 p
\textsc{السيّد بلانكو}. --- هوّي \VUgras{الدهّان} \textsuperscript{(102)}.
واقف على سلّمه مع علبة \VUgras{دهان} \textsuperscript{(103)}
و\VUgras{فرشايته} \textsuperscript{(104)}. عم يدهن خشب باب الدخول ليضلّ
أطول.

%%K t05-01-37 p
وهوّي كمان بيلزق الورق بالشقق.

%%K t05-01-38 p
و\VUgras{المنجّد} \textsuperscript{(107)} اللي على \VUgras{السلّم}
\textsuperscript{(106)}، عالـ\VUgras{بلكون} \textsuperscript{(28)}، عم
يحطّ \VUgras{ستارة} \textsuperscript{(29)}.

%%K t05-01-39 p
والعامل الواقف جنب الشبّاك الكبير هوّي \VUgras{عامل الإزاز}
\textsuperscript{(96)}. عم يثبّت \VUgras{ألواح الإزاز} \textsuperscript{(18)}
بـ\VUgras{المعجون} \textsuperscript{(73)}. وهداك العامل التاني، الراكع جنب
باب القبو، هوّي \VUgras{عامل الأقفال} \textsuperscript{(97)}. عم يركّب
\VUgras{قفل} \textsuperscript{(6)}. و\VUgras{مبرده} \textsuperscript{(98)}
جنبه.

%%K t05-01-40 p
\textsc{حنّا}. --- وكمان عامل أقفال هالعامل اللي عم يدقّ
بـ\VUgras{شاكوشه} \textsuperscript{(84)} على حديدة ؟

%%K t05-01-41 p
\textsc{السيّد بلانكو}. --- بالأصحّ، هوّي حدّاد (125). عم يطرق
\VUgras{حدايد} \textsuperscript{(67)} لـ\VUgras{الكهربجي} \textsuperscript{(124)}
اللي على سقف \VUgras{الكراج} \textsuperscript{(51)}. وعنده \VUgras{كور}
\textsuperscript{(126)} و\VUgras{سندان} \textsuperscript{(127)}
و\VUgras{كمّاشة} \textsuperscript{(128)}.

%%K t05-01-42 p
والكهربجي عم يثبّت \VUgras{سلك} التلفون \VUgras{النحاس}
\textsuperscript{(44)}.

%%K t05-tit-7 sub
\VUpk{10.2pt}{40}{شغل الجنينة.}
\VUsaut{3.0mm}

%%K t05-01-43 p
\textsc{العمّ}. --- بآخر \VUgras{الحوش} \textsuperscript{(45)}، جنب
\VUgras{الشبك} \textsuperscript{(47)} و\VUgras{البوّابة} \textsuperscript{(48)}
عم تشوف \VUgras{الجنايني} \textsuperscript{(129)}. عم يحضّر
\VUgras{الجنينة} \textsuperscript{(49)}. قلب التربة بـ\VUgras{مجرفته}
\textsuperscript{(131)}، وزرع البذور ومشّط بـ\VUgras{المشط}
\textsuperscript{(130)}. وبعدين، بـ\VUgras{مرشّته} \textsuperscript{(132)}،
رح يسقي \VUgras{الأحواض} \textsuperscript{(150)} والزرع رح يكبر.

%%K t05-01-44 p
و\VUgras{السايس} \textsuperscript{(133)} اللي هلّق اعتنى بالحصان وأكّله،
عم ينضّف \VUgras{العربايّة} \textsuperscript{(52)} بإسفنجة و\VUgras{مضخّة
اليد} \textsuperscript{(134)} وسطل ميّ. وبعدين رح يرجّعها عالكراج
ويسكّر الباب بـ\VUgras{الترباس} \textsuperscript{(53)} التخين.

%%K t05-01-45 p
هلّق مبسوط، بظنّ ؛ أخدت شرح كفاية.

%%K t05-01-46 p
\textsc{حنّا}. --- إي يا عمّي، بس كتير بحبّ شوف البيت من جوّا.

%%K t05-01-47 p
\textsc{العمّ}. --- هاد بكرا. السيّد روفو رح يشرحلك المخطّط ويقلّك اسم
كل غرفة.

%%K t05-tit-8 sub
\VUsaut{3.0mm}
\VUpk{10.2pt}{40}{ترتيب البيت من جوّا.}
\VUsaut{2.4mm}

%%K t05-apar-2 apar
{\centering\textit{(شوف المخطّط.)}\par}
\VUsaut{2.4mm}

%%K t05-01-48 p
\textsc{المهندس المعماري}. --- تعا يا حنّا، رح فرجيك حالاً أقسام
البيت.

%%K t05-01-49 p
خلّينا نفوت عالـ\VUgras{طابق الأرضي} \textsuperscript{(7)}. هيّو زرّ
\VUgras{الجرس} \textsuperscript{(11)}. منطلع \VUgras{الدرجات}
\textsuperscript{(14)} التلاتة تبع \VUgras{المصطبة} \textsuperscript{(9)}،
ومنعدّي \VUgras{عتبة} \textsuperscript{(10)} \VUgras{باب الدخول}
\textsuperscript{(8)}، وهيّنا صرنا عند أول \VUgras{الدرج}
\textsuperscript{(13)}.

%%K t05-01-50 p
\textsc{حنّا}. --- هاد الممرّ.

%%K t05-01-51 p
\textsc{المهندس المعماري}. --- لأ، هاد \VUgras{المدخل}
\textsuperscript{(12)}. \VUgras{الممرّ} \textsuperscript{(a)} بالآخر.
وعاليمين، هيّو \VUgras{الصالون} \textsuperscript{(b)} و\VUgras{غرفة
السفرة} \textsuperscript{(c)}، وقبال غرفة السفرة في \VUgras{المطبخ}
\textsuperscript{(d)} و\VUgras{غرفة الخدم} \textsuperscript{(e)}، وبعدين
لقدّام \VUgras{المكتب} \textsuperscript{(f)}. و\VUgras{الفرندة}
\textsuperscript{(23)} مع \VUgras{بابها الإزاز} \textsuperscript{(25)} جنب
الصالون.

%%K t05-01-52 p
هلّق منطلع الدرج. مسّك حالك بـ\VUgras{الدرابزين} \textsuperscript{(15)}
وعدّ الدرجات. هنّي أربعتعشر. هيّي \VUgras{الصالة} \textsuperscript{(m)}،
صرنا عالـ\VUgras{طابق} \textsuperscript{(27)} الأول.

%%K t05-tit-9 sub
\VUsaut{3.0mm}
\VUpk{10.2pt}{40}{الطابق الأول.}
\VUsaut{3.0mm}

%%K t05-01-53 p
قبال الدرج، في \VUgras{المرحاض} \textsuperscript{(20)} و\VUgras{غرفة
التواليت} \textsuperscript{(j)}. وعاليمين \VUgras{غرفة نوم}
\textsuperscript{(g)} حلوة مع شبّاكين على \VUgras{واجهة}
\textsuperscript{(1)} البيت. وعشمال، في \VUgras{الحمّام}
\textsuperscript{(1)} و\VUgras{غرفة الضيوف} \textsuperscript{(i)} مع
\VUgras{ليوان} \textsuperscript{(h)} كبير. وجنب غرفة النوم في \VUgras{غرفة
الولاد} \textsuperscript{(k)}. وهيّي عم تطلّ على \VUgras{تراس}
\textsuperscript{(21)} واسع. وتحت هالتراس في مرحاض تاني (20)، وعالحيط،
هيّو \VUgras{الجرس} \textsuperscript{(22)} اللي بيدقّ للعشا.

%%K t05-01-54 p
\textsc{حنّا}. --- خلّينا نطلع عالـ\VUgras{طابق التاني}.

%%K t05-01-55 p
\textsc{المهندس المعماري}. --- ما في طابق تاني. تحت السقف في بس
\VUgras{غرف السقف} \textsuperscript{(33)} والعلّية (34). خلّصنا الجولة،
فينا ننزل.

%%K t05-01-56 p
\textsc{حنّا}. --- شكراً يا سيّدي، هاد كتير مشوّق. حالاً رح قصّ لأبي كل
شي شفته.

\VUsaut{4.0mm}
\VUfilet{22mm}
